% ============================================
% Johns Hopkins Letter Template
% Assistant Professor Cover Letter – Villanova University
% ============================================

\documentclass[11pt,letterpaper]{letter}

\usepackage{graphicx}
\usepackage{eso-pic}
\usepackage{geometry}
\usepackage{microtype}
\usepackage[hidelinks]{hyperref}

% ----- Page Setup -----
\geometry{left=25mm,right=25mm,top=25mm,bottom=25mm}

% ----- Logo Placement (first page only) -----
\newcommand{\PutJHULogoFG}[1]{%
  \AddToShipoutPictureFG*{%
    \AtPageUpperLeft{%
      \hspace{10mm}%
      \raisebox{-38mm}{%
        \includegraphics[width=6cm]{#1}%
      }%
    }%
  }%
}

% ----- Sender Information -----
\signature{Decory Edwards}
\address{Department of Economics \\ Johns Hopkins University \\ Baltimore, MD 21218 \\ \texttt{dedwar65@jh.edu}}

\begin{document}
\PutJHULogoFG{Images/jhulogo.png}

\begin{letter}{Search Committee \\
Department of Economics \\
Villanova School of Business \\
Villanova University \\
800 Lancaster Avenue \\
Villanova, PA 19085 \\ USA}

\opening{Dear Members of the Search Committee,}

I am writing to express my interest in the tenure-track Assistant Professor position in the Department of Economics at Villanova University, beginning Fall 2026. I am a Ph.D. candidate in Economics at Johns Hopkins University, advised by Professors Christopher D.~Carroll, Nicholas W.~Papageorge, and Stelios Fourakis, and I expect to receive my degree in May 2026. My research fields are macroeconomics, wealth and income dynamics, and computational economics.

My research agenda focuses on understanding the determinants of wealth inequality and the mechanisms through which heterogeneity in returns to wealth shapes aggregate outcomes. In my job-market paper, I estimate the degree of heterogeneity in individual portfolio returns using micro-data from the Survey of Consumer Finances and simulate a structural model in which households face idiosyncratic return risk. I show that allowing for persistent heterogeneity in returns substantially improves the model’s ability to match the empirical distribution of wealth, offering a tractable way to connect micro-level return dispersion to macroeconomic inequality.

In a second project, I use the Health and Retirement Study to examine the relationship between interpersonal trust and portfolio performance. Building on the literature that finds a hump-shaped relationship between trust and income, I show that a similar relationship holds for individual returns to wealth measured in the HRS data, highlighting a behavioral channel through which social attitudes shape saving and investment outcomes. This work also provides new estimates of the persistent component of returns to net worth, which motivate the calibration of heterogeneity in my structural model.

Villanova University’s Department of Economics, situated within the Villanova School of Business, is an excellent fit for my research and teaching profile. The department’s search includes macroeconomics, financial economics, applied microeconomics, industrial organization, health economics, and law and economics—areas tightly connected to my work on heterogeneous returns, portfolio behavior, and inequality dynamics. I am particularly drawn to Villanova’s 2–2 teaching load, its strong support for high-quality publication, and its commitment to integrating rigorous applied research with business-relevant economic analysis. I also value Villanova’s Augustinian mission and its emphasis on community, service, and educating the whole person, which align with my own approach to teaching and mentorship.

\textbf{Teaching.} My teaching experience centers on undergraduate macroeconomics and an upper-level macro policy seminar, where I have led weekly discussion sections and held regular office hours for classes ranging from 16 to 40 students. My approach is student-centered: I aim to help students “convince themselves” that they have moved from confusion to understanding by holding structured office-hour coaching that builds trust and confidence. At Villanova, I am prepared to teach \textit{Principles of Macroeconomics}, \textit{Intermediate Macroeconomic Theory}, \textit{Money and Banking}, \textit{Econometrics}, and related business-oriented courses depending on departmental needs, at both the undergraduate and graduate levels.

I believe I would be a strong fit because my computational and empirical toolkit allows me to (i) produce rigorous, data-backed estimates of returns and distributional outcomes, (ii) build and validate quantitative models that speak to macroeconomic and business-economics questions, and (iii) mentor students effectively in a research-active environment. I am excited about the opportunity to contribute to Villanova University’s teaching, research, and service missions within the Villanova School of Business.

I would be honored to join the Department of Economics at Villanova University. Thank you for your consideration; I welcome the opportunity to discuss my fit with your department at your convenience.

\closing{Sincerely,}

\end{letter}
\end{document}
